\begin{claim}{F.18}{Proved}{The unique full branch lift and its exact source coset}
Keep the actual eight-point configuration $\mathcal P_8$
and involution $\mathsf r$ of F.12--17.  Put
$a_+=a(s_0)$ and $a_-=-a(s_0)$ with the literal ordering
in F.17.  Mark these two additional points and retain
the same homeomorphism as $\mathsf r^\dagger$.
Write
\[
\begin{split}
 P_{10}^{\dagger}
  &=\widehat{P\Gamma}(\mathbf P^1;\mathcal P_8,a_+,a_-),\\
 P_8&=\widehat{P\Gamma}(\mathbf P^1;\mathcal P_8),\\
 f:P_{10}^{\dagger}&\longrightarrow P_8,\qquad
 \mathcal J=\ker f,\\
 C^\dagger&=C_{P_{10}^{\dagger}}(\mathsf r^\dagger),
 \qquad C_8=C_{P_8}(\mathsf r).
\end{split}                                                     \tag{F.18.1}
\]
Here $f$ forgets exactly the two named added points.
The branch-permutation character
$\epsilon:C_8\twoheadrightarrow C_2$ is defined by
the actual decorated quotient of $\mathsf r$.
Then
\[
 f:C^\dagger\xrightarrow{\ \sim\ }C_8^+
       :=\ker\epsilon,\qquad
             C_{\mathcal J}(\mathsf r^\dagger)=1              \tag{F.18.2}
\]
are an isomorphism of profinite groups and its exact
kernel assertion.

For the action $R^\dagger$ on $\mathcal J$ by
$\mathsf r^\dagger$, the full continuous cocycles
\[
 Z^1(R^\dagger,\mathcal J)
     =\{\delta\in\mathcal J:\delta R^\dagger(\delta)=1\}
                                                               \tag{F.18.3}
\]
have precisely two $\mathcal J$-conjugacy classes.
Every order-two section has trivial centralizer in
$\mathcal J$.  In particular the map
\[
 b\longmapsto b^{-1}R^\dagger(b)
                                                               \tag{F.18.4}
\]
is a homeomorphism from $\mathcal J$ onto the
coboundary class.
More generally, for $k\in P_{10}^{\dagger}$ with
$f(k)\in C_8$, the class of
$k^{-1}R^\dagger(k)$ is the class labelled
$\epsilon(f(k))$.

The original residual group $\mathcal E$ of F.12
has a unique full centralizing branch lift
\[
 \ell:\mathcal E\longrightarrow C^\dagger,\qquad
                   f\ell=\widehat\phi|_{\mathcal E}.
                                                               \tag{F.18.5}
\]
It is continuous and injective.
This includes the proof that its branch character
is zero; no ordered drilling is assumed in advance.

Here is the exact remaining comparison to the
rational drilling of F.17.  Put
\[
\begin{split}
 L^\dagger&=\widehat{\pi_1(Y^\dagger,s_0)},\qquad
 L=\widehat{\pi_1(Y,s_0)},\\
 q:L^\dagger&\twoheadrightarrow L,\qquad
 \Phi^\dagger=\widehat{\phi^\dagger}_*,\\
 K^+&=\Phi^\dagger(\ker q)
          =\Phi^\dagger(L^\dagger)\cap\mathcal J .
\end{split}                                                     \tag{F.18.6}
\]
For $u\in\mathcal E$ and any $w\in q^{-1}(u)$,
define the actual element
\[
             b(w)=\ell(u)^{-1}\Phi^\dagger(w)\in\mathcal J .
                                                               \tag{F.18.7}
\]
It is the unique solution of
\[
 b(w)^{-1}R^\dagger(b(w))
       =\Phi^\dagger(w)^{-1}R^\dagger(\Phi^\dagger(w)).
                                                               \tag{F.18.8}
\]
The right $K^+$-orbit $b(w)K^+$ is independent of $w$
and defines a continuous map
\[
 \mathfrak o:\mathcal E\longrightarrow\mathcal J/K^+ .
                                                               \tag{F.18.9}
\]
No group structure on this coset space is assumed.
Its distinguished fibre is exactly
\[
\begin{split}
 \mathfrak o^{-1}(K^+)
   &=\{u\in\mathcal E:
                 \ell(u)\in\Phi^\dagger(L^\dagger)\}\\
   &=q\bigl((\Phi^\dagger)^{-1}(C^\dagger)\bigr).
\end{split}                                                     \tag{F.18.10}
\]
These are equalities for the specified full maps.
Whether this fibre is all of $\mathcal E$, and
whether $\mathcal E$ is trivial, remain unresolved.
\end{claim}

\begin{proof}
\emph{The two quotient normalizers and their actual inverse.}
On the eight-point sphere the involution is the
M\"obius map $r(V)=a_+^2/V$, with four nontrivial
marked pairs and unmarked fixed points $a_+,a_-$.
After marking them, the same map fixes each new
label and retains the four old pairs.
The decorated quotient before marking has four
ordinary orbit labels and two branch points;
after marking it has the same four orbit labels
and the two branch points ordered individually.

Use the discrete covering-normalizer sequence F.7.5
on both sides of the forgetting map.
Its kernel is the same deck group
$\langle\mathsf r\rangle=C_2$.
On either pure centralizer this kernel is trivial,
because its nonidentity element permutes the eight
old labels.  Descent therefore identifies the
discrete pure centralizer before drilling with a
finite-index subgroup $Q_0$ of the decorated quotient
mapping class group.  Its image after ordered
drilling is exactly the subgroup $Q_0^+$ preserving
the two branch labels individually.
Indeed, a representative in $Q_0^+$ has its unique
lift fixing all old labels, and that lift fixes the
two ordered ramification points as well.
Conversely every pure lift after drilling has those
two properties.  The inverse comparison is this
same unique pure lift.

The branch permutation on $Q_0$ is onto $C_2$.
To see this with the old labels retained, interchange
the two branch points along an arc avoiding the
four ordinary quotient labels.
Lift that quotient homeomorphism.
If its endpoint exchanges the two names in some
old pair, push that ordinary quotient point once
around one branch point and lift the isotopy from
the identity.  This changes exactly that pair
permutation, fixes the branch points individually,
and leaves all other pair names fixed.
Correcting each mismatch gives a pure upstairs
class still exchanging the two unmarked fixed
points.  Thus $[Q_0:Q_0^+]=2$.

The prime-order normalizer and full-topology
comparison F.7.4--5 applies to these two actual
involutions.  In particular both actual pure
centralizers are their discrete completions.
Their branch-permutation homomorphism is continuous
in the induced topology: the full topology admits
this finite $C_2$ quotient.  Its kernel is open;
the completion of $Q_0^+$ is exactly that kernel.
Completing the discrete inverse maps above gives
the first isomorphism in (F.18.2), for the
actual forgetting map $f$.  Intersecting its
kernel with $\mathcal J$ gives the second assertion.
No intersection of two unrelated closures is used.

\emph{The actual two-point kernel and its section classes.}
Write $S_8=\mathbf P^1\setminus\mathcal P_8$ and
\[
 J_{\rm d}=\pi_1(\operatorname{Conf}_2(S_8),(a_+,a_-)).
                                                               \tag{F.18.11}
\]
The action on this ordered configuration fibre is
$(z_+,z_-)\mapsto(rz_+,rz_-)$; it does not interchange
the two new labels.
The two successive Birman sequences, with the
completed comparison of F.14--15, identify
$\widehat{J_{\rm d}}$ with the actual subgroup
$\mathcal J$ in (F.18.1).
For completeness, projection to the first point
gives the split free-by-free extension
\[
              1\longrightarrow F_8
                \longrightarrow J_{\rm d}
                \longrightarrow F_7\longrightarrow1 .
                                                               \tag{F.18.12}
\]
A section is obtained by putting the second point
far out in a plane coordinate with one original
puncture at infinity, and then adjusting the
basepoint by a path in the fibre.
Characteristic finite-index subgroups of $F_8$
and the finite action of the free quotient on
their finite quotients give the full topology on
that kernel.  Completion is exact.
Mapping this sequence into the two completed
Birman sequences proves the asserted embedding
and equality with the forgetting kernel, just
as in F.15.11, with these actual ranks and labels.

Consider the discrete group
\[
            E_{\rm d}=J_{\rm d}\rtimes
                             \langle\mathsf r^\dagger\rangle .
                                                               \tag{F.18.13}
\]
It is the actual subgroup of the ten-point mapping
class group generated by the two-point kernel and
the indicated involution; the nontrivial coset is
detected by the old-label permutation.
Each order-two section has exactly those four
old pairs and fixes the two added labels.
Nielsen realization in the punctured form used
in F.15 realizes it as a sphere involution.
Its two fixed points are precisely the two
added marked points.

An ordered quotient homeomorphism matching those
two branch points and the four named ordinary
orbits lifts to a conjugator to
$\mathsf r^\dagger$.  Correct any old pair-name
mismatches by the quotient point pushes just
described.  This proves that every such section
is conjugate to $\mathsf r^\dagger$ by a pure
ten-point class $g$.
The condition that it lies in (F.18.13) is exactly
$f(g)\in C_{8,{\rm d}}$, the discrete centralizer.
The set of its $J_{\rm d}$-conjugacy classes is
therefore the actual double-coset set
\[
 J_{\rm d}\backslash
       \{g\in P\Gamma_{0,10}^{\dagger}:f(g)\in C_{8,{\rm d}}\}
       /C^\dagger_{\rm d}
       \ \xrightarrow{\ \sim\ }\
       C_{8,{\rm d}}/C_{8,{\rm d}}^+
       \ \cong C_2 .                                         \tag{F.18.14}
\]
The map sends $g$ to the branch permutation of $f(g)$.
It is onto because forgetting the two points is
onto on pure mapping class groups.
If two images have the same branch permutation,
their quotient downstairs lifts through
$C^\dagger_{\rm d}$ by the discrete isomorphism
already proved; removing that lift leaves an
element of $J_{\rm d}$.  This proves injectivity.
All their centralizers in $J_{\rm d}$ are trivial:
for the standard section this is the discrete
kernel of (F.18.2), and pure conjugation preserves
the normal subgroup $J_{\rm d}$.

Here are the full completion hypotheses for
transporting this calculation.
The group in (F.18.12) is good, by the two-row
Lyndon--Hochschild--Serre comparison for its free
kernel and free quotient, whose finite-coefficient
cohomology groups are finite.
It has a finite two-dimensional classifying model:
use finite graphs for the free kernel and quotient,
and mapping cylinders of graph homotopy equivalences
representing the finitely many free-group automorphisms.
The finite extension (F.18.13) is residually finite
as an actual subgroup of the residually finite
mapping class group.  Its completion is
$\mathcal J\rtimes\langle\mathsf r^\dagger\rangle$,
using the finite-index subgroup $J_{\rm d}$.
The finite-extension cohomology comparison gives
goodness and finite virtual cohomological type.
The same argument for all discrete two-primary
modules follows by taking the filtered union of
their finite invariant submodules; the
$\mathrm{FP}_\infty$ resolution commutes with
these unions.  These are precisely the
hypotheses used for (F.15.14).

Apply \cite[Theorem~A]{R-BZ} to this actual group
and the prime two.
It gives the two full conjugacy classes of
order-two subgroups and the closure of their
centralizers.  An $E_{\rm d}$-conjugacy between
sections is already a $J_{\rm d}$-conjugacy:
an outside conjugator can be multiplied by
the nonidentity element of the section it
conjugates to make its quotient label zero,
without changing that conjugation.
The same argument works after completion.
Thus (F.18.14) gives exactly two
$\mathcal J$-classes.
For each discrete section $T$, its centralizer
in $E_{\rm d}$ is $T$ itself, since its intersection
with $J_{\rm d}$ is trivial and the quotient is
$C_2$.  The completed centralizer equality therefore
makes the full fixed groups in $\mathcal J$
trivial as well.
The orbit map from $\mathcal J$ to either
cocycle class is consequently injective
and onto; compactness and the Hausdorff
target make it a homeomorphism.
This proves (F.18.3)--(F.18.4).
For a full $k$ as in the claim with
$\epsilon(f(k))=0$, let $c\in C^\dagger$
be the unique lift of $f(k)$ under (F.18.2).
Then $b=c^{-1}k\in\mathcal J$ and
$k^{-1}R^\dagger(k)=b^{-1}R^\dagger(b)$.
If instead $\epsilon(f(k))=1$, a coboundary
expression for this element would put
$kb^{-1}$ in $C^\dagger$, contradicting
(F.18.2) after forgetting.
Since there are exactly two full classes,
this proves the asserted class label for
every full $k$ directly.

\emph{The residual group fixes the two actual stabilizer labels.}
We verify this before defining the lift.
Let $\jmath:L\to\widehat P\Gamma_{0,12}$ be
the injective original covering map of F.12.
For $u\in\mathcal E$, both $\jmath(u)$ and
$\beta\jmath(u)\beta^{-1}$ lie in the pure
completed centralizer of $\tau$.
Their descents agree by the defining equation
of $\mathcal E$.  That descent is injective
on the pure centralizer: its only finite
kernel is $\langle\tau\rangle$, whose
nonidentity element has a nontrivial label
permutation.  Therefore
\[
                    \beta\jmath(u)\beta^{-1}=\jmath(u).
                                                               \tag{F.18.15}
\]
The element also centralizes $\tau$.
This establishes pointwise centralization
of $V_4=\langle\tau,\beta\rangle$; it is
stronger than preserving the two lifts
$\{\beta,\tau\beta\}$ as an unordered set.

The actual coordinate $W=(z-s_0)/(z+s_0)$
on the original sphere has the formulas
\[
 \tau(W)=-W,\qquad
 \beta(W)=\frac{k}{W},\qquad
 \tau\beta(W)=-\frac{k}{W},\qquad
 k=a(s_0)=\left(1-\frac{2\theta^2}{5}\right)(1+2i).
                                                               \tag{F.18.16}
\]
These follow by conjugating the original
M\"obius matrices of F.11--12; the exact
certificate below checks them in the
same coefficient field.
The quotient by $\tau$ is $V=W^2$.
The fixed points of $\beta$ consequently
give $V=k=a_+$, and those of $\tau\beta$
give $V=-k=a_-$.
Thus the literal F.17 branch ordering
is the ordering of these two distinct
stabilizer types.

A discrete class centralizing $V_4$
preserves each stabilizer subgroup
individually in the actual quotient
cover.  Its descent under $\tau$
therefore fixes these two labels
individually.
The full statement follows from the
direct prime-power comparison for the
actual two-group $V_4$: its completed
centralizer is the closure of its
discrete centralizer, with the full
topology of F.7.4--5.
The branch character is finite and
continuous on the target centralizer,
so the same zero character holds on
that entire closure.
Applying this to (F.18.15) proves
$\widehat\phi(\mathcal E)\subset C_8^+$.
The inverse in (F.18.2) now defines
the continuous injective map $\ell$
in (F.18.5), and proves its uniqueness.

\emph{The exact remaining source-image test.}
The square F.17.10 is the literal identity
\[
                       f\Phi^\dagger=\widehat\phi\,q .
                                                               \tag{F.18.17}
\]
Since $\widehat\phi$ is injective, an
element $\Phi^\dagger(w)$ lies in
$\mathcal J$ precisely when $q(w)=1$.
This proves both inclusions in the
equality defining $K^+$ in (F.18.6).
It is a closed subgroup, being the
compact image of the closed kernel
of $q$.

For $q(w)=u\in\mathcal E$, (F.18.17)
and (F.18.5) put $b(w)$ in $\mathcal J$.
Because $\ell(u)$ centralizes
$\mathsf r^\dagger$, direct multiplication
gives (F.18.8).
Equation (F.18.4) and its trivial fixed
kernel prove uniqueness of this solution.
Replacing $w$ by $wv$ with $v\in\ker q$
gives the exact formula
\[
                         b(wv)=b(w)\Phi^\dagger(v).
                                                               \tag{F.18.18}
\]
Thus $b(w)K^+$ depends only on $u$.
The map on $q^{-1}(\mathcal E)$ is
continuous; the compact surjection to
$\mathcal E$ is a quotient map.
It therefore descends continuously
to (F.18.9).

Finally,
$\ell(u)=\Phi^\dagger(w)b(w)^{-1}$.
The definition and equality in (F.18.6)
show that this element belongs to
$\Phi^\dagger(L^\dagger)$ precisely
when $b(w)\in K^+$.
In that case choose $v\in\ker q$
with $\Phi^\dagger(v)=b(w)^{-1}$;
then $\Phi^\dagger(wv)=\ell(u)$
centralizes and $q(wv)=u$.
Conversely, a source element with
centralizing ten-point image fills
to $\widehat\phi(u)$ with $u\in\mathcal E$,
and uniqueness in (F.18.2) identifies
its image with $\ell(u)$.
These are both directions of
(F.18.10).
In particular the canonical correcting
element is unique; its membership in
the specified closed subgroup $K^+$
is the remaining literal source test.

The executed certificate
\path{code/part_a/common_reflection_branch_lift_f18.py}
checks the original M\"obius conjugations,
the exact coefficient $k$, the four old
marked pairs, and the $V_4$ stabilizer
labels.  It passes 34 exact checks.
The full topological and completion
assertions are the arguments above,
not finite-quotient lifting assumptions.
\end{proof}
